\begin{center}\begin{tikzpicture}[font=\small,>=stealth,line join=round]
\draw[gray!40,rounded corners=3pt] (0,0) rectangle (14.2,8.5);\node at (7.1,7.9) {平移证明：把原来的对称轴搬到零点};
\node at (3.2,6.9){$f(x)=(x-2)^2+1$};\begin{scope}[shift={(1.2,2.4)},x=1cm,y=.62cm]\draw[->,gray] (-0.5,0)--(4.5,0);\draw[->,gray] (0,0)--(0,6);\draw[blue,thick,domain=0:4,samples=120] plot(\x,{(\x-2)^2+1});\draw[violet,dashed] (2,0)--(2,5.7);\node[below=6pt] at (2,0){$2$};\fill[orange!85!black] (2,1) circle(2pt);\end{scope}\node at (10.3,6.9){$g(t)=f(2+t)=t^2+1$};\begin{scope}[shift={(10.3,2.4)},x=1cm,y=.62cm]\draw[->,gray] (-2.5,0)--(2.5,0);\draw[->,gray] (0,0)--(0,6);\draw[blue,thick,domain=-2:2,samples=120] plot(\x,{(\x)^2+1});\draw[violet,dashed] (0,0)--(0,5.7);\node[below=6pt] at (0,0){$0$};\fill[orange!85!black] (0,1) circle(2pt);\end{scope}\draw[->,orange!85!black,thick] (6.1,4.3)--(7.2,4.3);\node at (7.1,1.5){原图向左平移2个单位；两图横纵单位分别相同，曲线高度不变。};
\node[text width=12.6cm,align=center] at (7.1,.65) {横坐标由 $x$ 改为 $t=x-2$，因此轴 $x=2$ 对应新的 $t=0$。};
\end{tikzpicture}\end{center}
